% Summary and Conclusion

\section{SUMMARY AND CONCLUSION}

\subsection{Summary}

This study designed, developed, and evaluated \textbf{VERACISM} (Verifiable Academic Credential
Integrity System with Multi-tenancy), a decentralized academic document verification platform
developed for the University of the Philippines Open University (UPOU). The system addresses the
persistent problem of academic credential fraud by combining cryptographic integrity verification,
decentralized Web3 storage, and multi layered threat scanning into a single cohesive platform.

VERACISM adopts a multitenant architecture that allows multiple educational institutions
(tenants) to independently manage their own documents while sharing common platform
infrastructure. The core workflow consists of three stages: (1) secure document upload with
real time malware and phishing scanning via ClamAV, YARA rules, and a PDF analyzer; (2) encrypted
decentralized storage on Filecoin/IPFS through the Lighthouse protocol, producing a tamper evident
Content Identifier (CID); and (3) public verification, where any stakeholder may confirm a
document's authenticity by resolving the CID and comparing its SHA-256 hash against the original
record.

The system was evaluated across five testing phases: functional and integration testing,
performance testing, security testing, and user acceptance testing.

All evaluation activities followed the baseline defined in the Project
Assessment chapter, including a hybrid test strategy that combined automated
API execution and manual role based verification.

\subsubsection{Functional and Integration Testing}

Automated testing executed a suite of \textbf{182 assertions} across 41 test cases: 22
integration test cases (93 assertions) covering the scan, upload, report retrieval, CID
verification, download, and storage workflows, and 19 functional test cases (89 assertions)
exercising role based access control for Admin and Tenant roles. All 182 assertions passed,
yielding a \textbf{100\% automated test pass rate} (see Appendix~\ref{appendix:integration-tests}
and Appendix~\ref{appendix:functional-tests}).

Role based manual and code review testing evaluated 30 test cases spanning Admin and Tenant
workflows. Of these, \textbf{27 passed fully} (90.0\%), \textbf{2 passed partially} (6.7\%), and
\textbf{1 failed} (3.3\%). The single failure (TC-ITO-005) identified the absence of a retry
policy when the Lighthouse storage service is temporarily unavailable, leaving the system
vulnerable to transient upload failures. Two partial results highlighted incomplete scoped audit
log visibility in the UI (TC-ITO-006) and the absence of a \texttt{VerificationLog} write during
public document verification (TC-E2E-003). These three defects have been recorded and are
addressed in the Future Work section. Despite these items, the core document lifecycle of upload,
scan, store, retrieve, and verify operated correctly across all evaluated scenarios.

\subsubsection{Performance Testing}

Performance evaluation was conducted in two rounds using k6 against the public CID verification
endpoint at 100 virtual users (VUs) sustained for five minutes.

The initial round (\textbf{Version 1.0}) produced an elevated error rate, which was subsequently
traced not to application layer defects but to an Nginx reverse proxy rate limiting policy that
throttled concurrent requests before they reached the backend. After the rate limiting threshold
was adjusted, the second round (\textbf{Version 1.1}) showed a \textbf{significantly reduced error
rate}, confirming that the application layer is performant under the tested load conditions. Both
rounds demonstrated that median response times remained within acceptable bounds, with P95 latency
values consistent with the expected overhead of decrypting and streaming content from the
Lighthouse proxy. A summary of both rounds is provided in
Appendix~\ref{appendix:performance}.

These outcomes satisfied the planned performance criteria in the Project
Assessment baseline, namely stable availability under concurrent traffic,
acceptable response behavior, and clear root cause identification for observed
errors.

\subsubsection{User Acceptance Testing}

User acceptance testing was completed using role based task scripts aligned to
the Project Assessment chapter. The executed UAT cycle covered 30 role based
cases across administrator, tenant office, registrar, student, external
verifier, and cross role flows.

UAT scoring was based on participant responses captured in
\textit{VERACISM\_UAT - Form Responses 1.csv}.

UAT was scored as role scoped responses, where non applicable UAT columns were
left blank by design for participants outside that role. Across 16 recorded
role scoped evaluations, outcomes were \textbf{15 passed} (93.8\%),
\textbf{1 partially completed} (6.2\%), and \textbf{0 failed}. At the
unique UAT case level, all 8 UAT IDs received responses, with 7 passed and 1
partial outcome.

These outcomes confirm end user readiness for controlled deployment while
preserving visibility of deferred improvements already captured in the
roadmap (see Appendix~\ref{appendix:uat-results}).

\subsubsection{Security Testing}

Security assessment was performed in two phases. The first phase employed Burp Suite
Community Edition as a Dynamic Application Security Testing (DAST) tool and produced several
tentative findings including SQL injection and reflected XSS candidates. Upon investigation, all
SQL injection findings were determined to be false positives: the backend exclusively uses Entity
Framework Core with parameterized queries, eliminating the SQL injection attack surface.
Reflected XSS candidates were similarly reviewed against the Angular frontend, which enforces
strict contextual output encoding by default.

The second phase used manual security testing and code level remediation review to validate
the initial scanner output. This phase confirmed that many findings from the first DAST pass
were false positives rather than exploitable vulnerabilities. After validation, the team identified
and confirmed
\textbf{five genuine vulnerabilities} (VUL-01 through VUL-05) ranging in severity from Low to
Medium-High. All five were remediated prior to the final evaluation:

\begin{itemize}
    \item \textbf{VUL-01} (Low): Verbose error messages exposing internal stack traces;
          remediated by replacing raw exception output with generic error responses.
    \item \textbf{VUL-02} (Medium-High): Missing \texttt{INTERNAL\_API\_KEY} enforcement on the
          Lighthouse encryption service; remediated by implementing API key middleware.
    \item \textbf{VUL-03} (Low-Medium): Multer-uploaded temporary files not deleted after
          processing; remediated by adding deterministic cleanup logic.
    \item \textbf{VUL-04} (Medium-High): Missing \texttt{GET /health} endpoint causing Docker
          healthcheck failures; remediated by adding a health endpoint.
    \item \textbf{VUL-05} (Medium-High): Inline \texttt{HttpClient} instantiation bypassing
          the \texttt{LighthouseStorageService} abstraction; remediated by routing all Lighthouse
          calls through the service layer.
\end{itemize}

Detailed vulnerability evidence and remediation steps are provided in
Appendix~\ref{appendix:security}.

The security outcomes were also consistent with the Project Assessment
baseline, which required verification of access boundaries, safe error
handling, and evidence based review of scanner findings.



% -----------------------------------------------------------------------
\subsection{Conclusion}

The results collectively demonstrate that VERACISM is a \textbf{operationally validated, security hardened,
and reasonably performant} platform for decentralized academic credential verification. The system
successfully fulfills its primary objective: enabling educational institutions to upload student
documents, secure them with cryptographic integrity proofs on a decentralized network, and allow
any stakeholder to independently verify authenticity without requiring access to a central
authority or trusted intermediary.

The automated test suite recorded a \textbf{100\% assertion pass rate}, and the manual test suite
achieved a \textbf{90.0\% full pass rate} with identified gaps limited to noncritical features
(retry logic, audit log UI scoping, and verification log persistence). These gaps do not compromise
the core verification guarantee and are scoped as incremental improvements in the roadmap.

The security posture of VERACISM improved substantially over the course of development. All five
confirmed vulnerabilities were remediated, and the DAST flagged SQL injection findings were
eliminated by architectural choice: the consistent use of Entity Framework Core with parameterized
queries throughout the backend. Post remediation, the system presents no known critical or high
severity vulnerabilities.

Performance testing confirmed that the application layer is capable of handling moderate concurrent
load (100 VUs) without service degradation, provided that infrastructure level rate limiting is
tuned appropriately to the expected request volume. The discovery of the Nginx rate limiting
misconfiguration during Version 1.0 testing, and its resolution in Version 1.1, is itself a
valuable outcome: it demonstrates that performance testing surfaced a real infrastructure gap that
would not have been visible through functional testing alone.

\subsubsection{Contributions}

This study contributes: (1) a working open source reference implementation of a
multitenant decentralized document verification system using Web3 storage; (2) a DevSecOps
integration pattern that embeds ClamAV, YARA, and PDF phishing analysis into the document
ingestion pipeline; (3) empirical evidence from functional, integration, performance, and
security testing of the platform; and (4) a documented vulnerability remediation lifecycle
applicable to similar academic information systems developed in constrained institutional
environments.

VERACISM establishes a foundation for further research into blockchain anchored academic
credential ecosystems, privacy preserving verification protocols, and standards based
interoperability with frameworks such as W3C Verifiable Credentials and Open Badges.


